
%----------------------------------------------------------------------------------------
%	PACKAGES AND OTHER DOCUMENT CONFIGURATIONS
%----------------------------------------------------------------------------------------

\documentclass{resume} % Use the custom resume.cls style

\usepackage[UTF8,fontset=fandol]{ctex}
\usepackage[hidelinks]{hyperref}
\hypersetup{pdftitle={Ruochen Dai - Curriculum Vitae},pdfauthor={Ruochen Dai}}
\usepackage[left=0.75in,top=0.6in,right=0.75in,bottom=0.6in]{geometry} % Document margins
\newcommand{\tab}[1]{\hspace{.2667\textwidth}\rlap{#1}}
\newcommand{\itab}[1]{\hspace{0em}\rlap{#1}}
\linespread{1.0}
\setlength{\parindent}{0pt}
\name{Ruochen Dai} % Your name
\address{Central University of Finance and Economics}
% Your address
\address{ Shunsha Road, Changping District, Beijing} % Your secondary addess (optional)
\address{dairuochenpku@gmail.com} % Public email only

\begin{document}
\begin{center}Updated October 2026\end{center}
\pagestyle{plain}

%----------------------------------------------------------------------------------------
%	EDUCATION SECTION
%----------------------------------------------------------------------------------------

\begin{rSection}{Work Experience}

{\bf Central University of Finance and Economics, Beijing} \hfill {\em August 2019 - Present} 
\\ Associate Professor (with tenure)
\\ School of Economics

{\bf Peking University, Beijing} \hfill {\em July 2017 - June 2019} 
\\ Postdoctoral Research Associate 
\\ National School of Development
\end{rSection}

\begin{rSection}{Education}

{\bf Peking University, Beijing} \hfill {\em September 2013 - July 2017} 
\\ Ph.D. in Economics
\\ National School of Development

{\bf Peking University, Beijing} \hfill {\em September 2010 - July 2012} 
\\ B.S. in Statistics
\\ School of Mathematical Science

{\bf Peking University, Beijing} \hfill {\em September 2009 - July 2013} 
\\ B.A. in Economics
\\ School of Economics
\end{rSection}

\begin{rSection}{Research Interests}
{\bf Firm Dynamics in Developing Countries, Entrepreneurship, Development Economics} 
\end{rSection}

%--------------------------------------------------------------------------------
%    Projects And Seminars
%-----------------------------------------------------------------------------------------------



\begin{rSection}{Professional Activities}

\begin{rSubsection}{Field survey}{}{}{}
\item Enterprise Survey for Innovation and Entrepreneurship in China. 2017-2025
\item Survey of E-commerce Cluster in Baigou. 2015
\end{rSubsection}

\begin{rSubsection}{Referee Service}{}{}{}
\item {\em American Economic Journal: Macro}, {\em American Economic Journal: Economic Policy}, {\em Journal of Development Economics} , {\em Journal of Urban Economics} , {\em China Economic Review} , {\em Journal of Regional Science}, {\em Economic Development and Culture Change} 
\end{rSubsection}



\end{rSection}




\clearpage
\begin{rSection}{English Publications}
\begin{enumerate}
\item Loren Brandt, \textbf{Ruochen Dai}, Gueorgui Kambourov, Kjetil Storesletten, Xiaobo Zhang. (2026). \href{https://ruochendai-econ.github.io/papers/serial-entrepreneurship.pdf}{Serial Entrepreneurship in China}. \emph{American Economic Journal: Macroeconomics}, forthcoming.

\item \textbf{Ruochen Dai}, Xiaobo Zhang. (2026). \href{https://ruochendai-econ.github.io/papers/ecommerce-entrepreneurship.pdf}{E-commerce Expands the Bandwidth of Entrepreneurship}. \emph{The Journal of Law, Economics, and Organization}, 42(3), 906--933.

\item \textbf{Ruochen Dai}, Yue Feng, Chengfang Liu, Langrui Li, Lina Zhang, Ketong Zhu. (2026). \href{https://ruochendai-econ.github.io/papers/sme-formal-employment.pdf}{Why Do Chinese SMEs Avoid Formal Employment? Political Connections and Unequal Access to Formalization Benefits}. \emph{China Economic Review}, 95, 102638.

\item \textbf{Ruochen Dai}, Dongmei Guo, Yajie Han, Yu Qin. (2024). \href{https://ruochendai-econ.github.io/papers/pollution-induced-trips.pdf}{Pollution-Induced Trips: Evidence from Flight and Train Bookings in China}. \emph{Journal of Development Economics}, 171, 103340.

\item \textbf{Ruochen Dai}, Dilip Mookherjee, Yingyue Quan, Xiaobo Zhang. (2021). \href{https://ruochendai-econ.github.io/papers/clusters-covid.pdf}{Industrial Clusters, Networks and Resilience to the Covid-19 Shock in China}. \emph{Journal of Economic Behavior \& Organization}, 183, 433--455.

\item \textbf{Ruochen Dai}, Hao Feng, Junpeng Hu, Quan Jin, Huiwen Li, Ranran Wang, Ruixin Wang, Lihe Xu, Xiaobo Zhang. (2021). \href{https://ruochendai-econ.github.io/papers/smes-covid.pdf}{The Impact of COVID-19 on Small and Medium-Sized Enterprises (SMEs): Evidence from Two-Wave Phone Surveys in China}. \emph{China Economic Review}, 67, 101607.
\end{enumerate}
\end{rSection}

\begin{rSection}{Research Grants}
\begin{enumerate}
\item \textbf{General Program, National Natural Science Foundation of China}. A Unified National Market and the Spatial Agglomeration of Innovation: A Quantitative Study Using an Endogenous Growth Model with Knowledge Spillovers.\\ {\small 全国统一大市场与创新空间集聚------基于包含知识溢出内生增长模型的量化研究}\\ Grant No. 72673217; January 2027--December 2030.

\item \textbf{Young Scientists Fund, National Natural Science Foundation of China}. Digital Economy, Entrepreneurship and Innovation, and Modern Industrial Clusters.\\ {\small 数字经济，创业创新与现代化产业集群}\\ Grant No. 72203252; January 2023--December 2025.
\end{enumerate}\end{rSection}
\clearpage
\begin{rSection}{Publications in Chinese}
\begin{enumerate}
\item 陈斌开, 吴楠, \textbf{戴若尘}. (2025). \href{https://ruochendai-econ.github.io/papers/innovation-policy.pdf}{引入技术外溢视角的中国创新政策}. \emph{中国社会科学}(12), 69--87.

\item 冯悦, \textbf{戴若尘}, 陈斌开, 智镜伊. (2025). \href{https://ruochendai-econ.github.io/papers/skills-entrepreneurship.pdf}{先前职业技能、营商环境与企业家精神}. \emph{管理科学学报}, 28(11), 79--93.

\item \textbf{戴若尘}, 索昊, 陈斌开. (2025). \href{https://ruochendai-econ.github.io/papers/factor-market-integration.pdf}{要素市场一体化的测度与影响：微观企业数据的新视角}. \emph{世界经济}, 48(7), 105--132.

\item 赵扶扬, 王都, \textbf{戴若尘}, 陈斌开. (2025). \href{https://ruochendai-econ.github.io/papers/spatial-capital-allocation.pdf}{跨地投资扭曲、产品市场分割与资本的空间配置效率------基于量化空间均衡模型的研究}. \emph{管理世界}, 41(6), 20--43.

\item \textbf{戴若尘}, 冯悦, 丁睿, 马雪俨, 梁延东. (2024). \href{https://ruochendai-econ.github.io/papers/firm-digitalization.pdf}{中国企业数字化进程的典型事实------来自中国企业创新创业调查的证据}. \emph{经济管理学刊}, 3(2), 55--82.

\item \textbf{戴若尘}, 祝仲坤, 张晓波. (2024). \href{https://ruochendai-econ.github.io/papers/regional-entrepreneurship-index.pdf}{中国区域创新创业指数构建与空间格局（1990---2020）}. \emph{经济科学}(1), 5--34.

\item \textbf{戴若尘}, 王艾昭, 陈斌开. (2022). \href{https://ruochendai-econ.github.io/papers/digital-economy-index.pdf}{中国数字经济核心产业创新创业：典型事实与指数编制}. \emph{经济学动态}(4), 29--48.

\item 张睿, 张勋, \textbf{戴若尘}. (2018). \href{https://ruochendai-econ.github.io/papers/infrastructure-productivity.pdf}{基础设施与企业生产率：市场扩张与外资竞争的视角}. \emph{管理世界}, 34(1), 88--102.
\end{enumerate}
\end{rSection}

\begin{rSection}{Working Papers}
\begin{enumerate}
\item \textbf{Ruochen Dai}, Dilip Mookherjee, Kaivan Munshi, Xiaobo Zhang. (2026). \href{https://ruochendai-econ.github.io/papers/community-networks.pdf}{Community Networks and Entrepreneurship in a Developing Economy}. 

\item \textbf{Ruochen Dai}, Kailin Gao, Rui Zhang. (2026). \href{https://ruochendai-econ.github.io/papers/export-automation.pdf}{Export Destinations, Quality, and Automation: Evidence from China}. 
\end{enumerate}
\end{rSection}

\end{document}
